\subsection{One algorithm for all slices}\label{exch:sec-5}

Take \(T\ge1\); the case \(T=0\) is immediate. For \(\epsilon\in(0,1)\)
choose

\[
a=(\epsilon/(64T))^2,\qquad s=\max\{1,S(T,\epsilon/4,a)\},
\] \[
\Gamma=\min\{\epsilon/(8s),(\epsilon/(32T))^2\},\qquad
\eta=a\Gamma/2,
\] \[
L=\lceil4/(\epsilon\eta^2)\rceil,\qquad
N_0=\lceil1/(a^2\Gamma^2)\rceil,\qquad
M=\max\{2L,L+N_0\}.
\]

Each quantity is polynomially bounded in \(T,1/\epsilon\). If \(n<M\),
read every bit and output \(Q(X)\) exactly.

Otherwise query the first \(L\) bits, calculate \(q\), and restrict
\(Q\) to this observed prefix. Write the restricted acceptance function
as \(Q_{\rm tr}\) on the remaining \(N=n-L\) bits. Fixing known
coordinates does not increase quantum query complexity.

\begin{itemize}
\tightlist
\item
  If \(q\le2a\), output \(Q_{\rm tr}(0^N)\).
\item
  If \(q\ge1-2a\), output \(Q_{\rm tr}(1^N)\).
\item
  Otherwise run the product-bias-\(q\) simulator for \(Q_{\rm tr}\) from
  Section~\ref{exch:sec-2}, with error \(\epsilon/4\)
  and at most \(s\) further queries.
\end{itemize}

A suitable tree can be hardwired for each of the finitely many observed
prefixes. This uses unrestricted computation but at most \(L+s\)
queries, and never uses the unknown total weight \(k\).
